\textbf{COLUNI 2023 -- Questão 09}\par\medskip Leia o trecho a seguir.

\par

Conseguir o Clocky naturalmente não seria fácil, \underline{mas}, com a ajuda de amigos que estudavam no Instituto Tecnológico de Massachusetts, obteve uma cópia do projeto, com a promessa de mantê-lo em segredo. Já confeccionar o despertador não foi difícil: ele era técnico em eletrônica e tinha uma habilidade incrível. \underline{Assim}, logo tinha em sua mesa de cabeceira a engenhoca. […] Uma manhã, \underline{contudo}, Clocky ultrapassou todos os limites. Tocava como um demônio, e ia de peça em peça, seu dono correndo atrás. Clocky, o implacável. SCLIAR, Moacyr. \textbf{Histórias que os jornais não contam}. 3. ed. Porto Alegre: L\&PM Editores, 2017. p. 129-130.(Adaptado).

\par

As palavras destacadas no trecho acima: mas, Assim, e contudo expressam, respectivamente, a ideia de:

\par

a) contraste, conclusão e adversidade.                               c) comparação, conclusão e contraste. b) adversidade, explicação, conclusão.                               d) explicação, consequência e adversidade.
